\chapter{La rutina diaria con los scripts}\label{cap:rutina}

\section{Antes de la sesión (30--45 min)}
\begin{enumerate}
\item \textbf{Calendario}: datos macro con hora y consenso, bancos centrales, resultados de empresas,
  inventarios (EIA), vencimientos de opciones y futuros.
\item \textbf{Contexto}: qué hicieron Asia y Europa, futuros de índices, dólar, bonos a 10 años, crudo,
  oro, VIX. ¿Régimen de riesgo (\emph{risk-on}) o de aversión (\emph{risk-off})?
\item \textbf{Hoja premercado}: niveles, volatilidad y régimen de tu lista de seguimiento.
\begin{lstlisting}
python scripts/premercado.py --csv datos/SPY.csv --csv datos/EURUSD.csv --csv datos/CL.csv
\end{lstlisting}
  Imprime el checklist, las sesiones FX abiertas, ATR, volatilidad, RSI, distancia a las medias,
  pivotes, máximos/mínimos de 20 sesiones y el rango esperado del día.
\item \textbf{Escenarios si/entonces}: «si rompe 1,0900 con volumen tras el IPC, compro con stop 1,0870».
\item \textbf{Riesgo}: revisa el riesgo abierto y la cartera.
\begin{lstlisting}
python scripts/informe_riesgo.py --csv datos/cartera.csv --pos SPY=50000 --pos TLT=-20000 --pos GLD=15000
\end{lstlisting}
\end{enumerate}

\section{Durante la sesión}
\begin{itemize}
\item Calcula el tamaño \emph{antes} de enviar la orden:
\begin{lstlisting}
python scripts/tamano_posicion.py futuro --capital 250000 --riesgo 0.75 --simbolo CL --entrada 74.2 --stop 72.6
python scripts/tamano_posicion.py fx --capital 20000 --par USDJPY --tipo 151.2 --stop-pips 40
\end{lstlisting}
\item Envía stop y objetivo con la entrada (orden bracket).
\item Calculadoras rápidas: opciones, IV, movimiento esperado, forward, cruces, curva, crack, bonos, funding
  (\py{scripts/calculadoras.py --help}).
\item Respeta el límite de pérdida diaria.
\end{itemize}

\section{Después de la sesión}
\begin{enumerate}
\item Anota cada operación en el diario (\py{scripts/diario.py nueva}).
\item Puntúa el cumplimiento del plan del día.
\item Semanalmente: \py{scripts/diario.py resumen}, revisión de setups y de errores.
\item Mensualmente: re-evalúa estrategias sistemáticas con \py{scripts/backtest.py} y compara el
  resultado real con el esperado.
\end{enumerate}

\section{Datos}
Los scripts leen CSV con columna de fecha y \texttt{open, high, low, close[, volume]} (o una columna por
activo para el informe de riesgo). Puedes exportarlos desde tu bróker o plataforma; si instalas
\texttt{yfinance}, \py{trader.data.download("SPY")} descarga históricos diarios. Sin
\texttt{--csv}, todo funciona con datos sintéticos para practicar.

\section{Mapa de scripts}
\begin{table}[h]\centering\small
\begin{tabularx}{\textwidth}{lX}
\toprule
Script & Para qué \\ \midrule
\py{premercado.py} & Checklist, sesiones FX, niveles y régimen de cada activo. \\
\py{tamano\_posicion.py} & Tamaño por riesgo: acciones, FX, futuros, ATR, cripto (con aviso de liquidación). \\
\py{informe\_riesgo.py} & Vol, VaR/ES, contribuciones, correlaciones y estrés de la cartera. \\
\py{backtest.py} & Backtest de SMA, Donchian, RSI y pares con costes, métricas y gráfico. \\
\py{diario.py} & Diario de operaciones en CSV y estadísticas en $R$ por setup. \\
\py{calculadoras.py} & Opciones, IV, movimiento esperado, forward FX, cruces, curva, crack, bonos, funding. \\
\py{generar\_figuras.py} & Regenera las figuras de este PDF. \\
\bottomrule
\end{tabularx}
\end{table}
